% Môn: Vật lý
% Lớp: 12
% Chương: Chương I. Vật lí nhiệt
% Tên đề: Ôn tập Chương 1 - Lý 12 - Đề ôn 1 - Tân Bình
% Loại: Trắc nghiệm 1 phương án
% Trường: THPT Tân Bình

\section{ÔN TẬP CHƯƠNG I. VẬT LÍ NHIỆT}
\subsection{Đề ôn 1 -- Tân Bình}

\begin{ex}
Chọn câu \textbf{sai} khi nói về cấu tạo chất.
\choice
{Các phân tử luôn luôn chuyển động không ngừng.}
{Các phân tử chuyển động càng nhanh thì nhiệt độ của vật càng cao và ngược lại.}
{\True Các phân tử luôn luôn đứng yên và chỉ chuyển động khi nhiệt độ của vật càng cao.}
{Các chất được cấu tạo từ các hạt riêng biệt là phân tử.}
\loigiai{Các phân tử cấu tạo nên vật luôn chuyển động không ngừng. Vì vậy phát biểu đã cho là sai.}
\end{ex}

\begin{ex}
Chỉ ra kết luận \textbf{sai} trong các kết luận sau.
\choice
{Sự chuyển từ thể rắn sang thể lỏng gọi là sự nóng chảy.}
{Sự chuyển từ thể lỏng sang thể rắn gọi là sự đông đặc.}
{Trong thời gian nóng chảy (hay đông đặc) nhiệt độ của hầu hết các vật không thay đổi.}
{\True Các kim loại có nhiệt độ nóng chảy giống nhau.}
\loigiai{Mỗi chất có nhiệt độ nóng chảy riêng nên các kim loại khác nhau thường có nhiệt độ nóng chảy khác nhau.}
\end{ex}

\begin{ex}
Điều nào sau đây đúng khi nói về cấu tạo chất của thể khí?
\choice
{Khoảng cách giữa các phân tử gần nhau.}
{Sự sắp xếp của các phân tử có trật tự.}
{\True Các phân tử chuyển động hỗn loạn.}
{Các phân tử chỉ dao động quanh vị trí cân bằng cố định.}
\loigiai{Ở thể khí, các phân tử ở cách xa nhau và chuyển động hỗn loạn không ngừng.}
\end{ex}

\begin{ex}
Một quả bóng khối lượng $100\,\mathrm{g}$ rơi từ độ cao $10\,\mathrm{m}$ xuống sân và nảy lên được $7\,\mathrm{m}$. Lấy $g=9{,}8\,\mathrm{m/s^2}$. Độ biến thiên nội năng của quả bóng trong quá trình trên bằng
\choice
{\True $2{,}94\,\mathrm{J}$}
{$3{,}00\,\mathrm{J}$}
{$294\,\mathrm{J}$}
{$6{,}86\,\mathrm{J}$}
\loigiai{Độ giảm cơ năng chuyển thành nội năng: $\Delta U=mg(h_1-h_2)=0{,}1\cdot9{,}8\cdot(10-7)=2{,}94\,\mathrm{J}$.}
\end{ex}

\begin{ex}
Nội dung nguyên lí I nhiệt động lực học là:
\choice
{\True Độ biến thiên nội năng của vật bằng tổng công và nhiệt lượng mà vật nhận được.}
{Độ biến thiên nội năng của vật bằng tổng nhiệt lượng mà vật nhận được.}
{Độ biến thiên nội năng của vật bằng tổng công mà vật nhận được.}
{Độ biến thiên nội năng của vật bằng hiệu số công và nhiệt lượng mà vật nhận được.}
\loigiai{Theo nguyên lí I nhiệt động lực học: $\Delta U=A+Q$.}
\end{ex}

\begin{ex}
Trường hợp nội năng của vật bị biến đổi \textbf{không phải do truyền nhiệt} là:
\choice
{Chậu nước để ngoài nắng một lúc nóng lên.}
{Gió mùa đông bắc tràn về làm cho không khí lạnh đi.}
{\True Khi trời lạnh, ta xoa hai bàn tay vào nhau cho ấm lên.}
{Cho cơm nóng vào bát thì bưng bát cũng thấy nóng.}
\loigiai{Xoa hai bàn tay vào nhau làm nội năng tăng do thực hiện công.}
\end{ex}

\begin{ex}
Hệ thức $\Delta U=A+Q$ khi $Q<0$ và $A>0$ mô tả quá trình
\choice
{hệ truyền nhiệt và sinh công.}
{hệ nhận nhiệt và sinh công.}
{\True hệ truyền nhiệt và nhận công.}
{hệ nhận nhiệt và nhận công.}
\loigiai{$Q<0$ cho biết hệ truyền nhiệt ra môi trường; $A>0$ cho biết hệ nhận công.}
\end{ex}

\begin{ex}
Mỗi độ chia ($1^\circ\mathrm{C}$) trong thang Celsius bằng $X$ của khoảng cách giữa nhiệt độ tan chảy của nước tinh khiết đóng băng và nhiệt độ sôi của nước tinh khiết ở áp suất tiêu chuẩn. $X$ là
\choice
{$\dfrac{1}{273{,}16}$}
{\True $\dfrac{1}{100}$}
{$\dfrac{180}{273{,}15}$}
{$\dfrac{1}{273{,}15}$}
\loigiai{Khoảng từ $0^\circ\mathrm{C}$ đến $100^\circ\mathrm{C}$ được chia thành $100$ phần bằng nhau nên $X=\dfrac{1}{100}$.}
\end{ex}

\begin{ex}
Điểm đóng băng và sôi của nước theo thang Fahrenheit ở áp suất tiêu chuẩn là
\choice
{$0^\circ\mathrm{F}$ và $100^\circ\mathrm{F}$}
{$100^\circ\mathrm{F}$ và $200^\circ\mathrm{F}$}
{\True $32^\circ\mathrm{F}$ và $212^\circ\mathrm{F}$}
{$20^\circ\mathrm{F}$ và $220^\circ\mathrm{F}$}
\loigiai{$0^\circ\mathrm{C}=32^\circ\mathrm{F}$ và $100^\circ\mathrm{C}=212^\circ\mathrm{F}$.}
\end{ex}

\begin{ex}
Nhiệt độ trung bình ngày--đêm tại Hà Nội là $24^\circ\mathrm{C}$--$17^\circ\mathrm{C}$. Sự chênh lệch nhiệt độ này trong thang Kelvin là bao nhiêu?
\choice
{\True $7\,\mathrm{K}$}
{$24\,\mathrm{K}$}
{$17\,\mathrm{K}$}
{$280\,\mathrm{K}$}
\loigiai{$\Delta T=24-17=7\,\mathrm{K}$. Độ chênh lệch nhiệt độ trong Celsius và Kelvin có cùng độ lớn.}
\end{ex}

\begin{ex}
Trong thang nhiệt độ Kelvin, nhiệt độ của nước đá đang tan là $273\,\mathrm{K}$. Hỏi nhiệt độ của nước đang sôi là bao nhiêu?
\choice
{$0\,\mathrm{K}$}
{\True $373\,\mathrm{K}$}
{$173\,\mathrm{K}$}
{$100\,\mathrm{K}$}
\loigiai{$T_{\mathrm{sôi}}=273+100=373\,\mathrm{K}$.}
\end{ex}

\begin{ex}
Công thức tính nhiệt lượng cần truyền cho vật có khối lượng $m$ để làm vật nóng lên thêm $\Delta T$ độ K là
\choice
{\True $Q=mc\Delta T$}
{$Q=m\Delta T$}
{$Q=c\Delta T$}
{$Q=mcT$}
\loigiai{Nhiệt lượng cần truyền cho vật: $Q=mc\Delta T$.}
\end{ex}

\begin{ex}
Nhiệt hóa hơi riêng là
\choice
{lượng nhiệt cần thiết để chuyển $1\,\mathrm{kg}$ chất từ thể rắn sang thể lỏng.}
{\True lượng nhiệt cần thiết để chuyển $1\,\mathrm{kg}$ chất từ thể lỏng sang thể khí.}
{lượng nhiệt cần thiết để chuyển $1\,\mathrm{kg}$ chất từ thể khí sang thể lỏng.}
{lượng nhiệt cần thiết để chuyển $1\,\mathrm{kg}$ chất từ thể lỏng sang thể rắn.}
\loigiai{Nhiệt hóa hơi riêng là nhiệt lượng cần cung cấp để hóa hơi hoàn toàn $1\,\mathrm{kg}$ chất lỏng ở nhiệt độ sôi.}
\end{ex}

\begin{ex}
Quá trình biến đổi trạng thái của một khối nước đá từ trạng thái rắn sang trạng thái lỏng được biểu diễn bằng đồ thị nhiệt độ theo thời gian. Đoạn đồ thị ứng với quá trình nóng chảy của khối nước đá là
\choice
{đoạn $AB$.}
{\True đoạn $BC$.}
{đoạn $CD$.}
{đoạn $AB$ và $CD$.}
\loigiai{Trong quá trình nóng chảy của nước đá tinh khiết, nhiệt độ không đổi. Đoạn nằm ngang $BC$ biểu diễn quá trình nóng chảy.}
\end{ex}

\begin{ex}
Trong thí nghiệm đo nhiệt hóa hơi riêng của nước, nhiệt hóa hơi riêng của nước được xác định bằng công thức
\[
L=\frac{\mathcal{P}\tau}{\Delta m}.
\]
Giá trị $\Delta m$ là khối lượng
\choice
{nước ban đầu trong ấm đun.}
{nước trong ấm đun tại thời điểm $\tau$.}
{\True nước bị bay hơi sau thời gian $\tau$.}
{nước và ấm đun.}
\loigiai{Ta có $\mathcal{P}\tau=L\Delta m$, nên $\Delta m$ là khối lượng nước bị bay hơi sau thời gian $\tau$.}
\end{ex}

\begin{ex}
Biết nhiệt nóng chảy riêng của thiếc là $\lambda=0{,}61\times10^5\,\mathrm{J/kg}$. Nhiệt lượng cần cung cấp cho một cuộn thiếc hàn có khối lượng $0{,}4\,\mathrm{kg}$ nóng chảy hoàn toàn ở nhiệt độ nóng chảy là
\choice
{\True $24400\,\mathrm{J}$}
{$4880\,\mathrm{J}$}
{$4{,}88\times10^7\,\mathrm{J}$}
{$76250\,\mathrm{J}$}
\loigiai{$Q=\lambda m=0{,}61\times10^5\cdot0{,}4=24400\,\mathrm{J}$.}
\end{ex}

\begin{ex}
Tính nhiệt lượng cần cung cấp cho $4\,\mathrm{kg}$ nước đá ở $0^\circ\mathrm{C}$ để chuyển nó thành nước ở $20^\circ\mathrm{C}$. Biết $\lambda=3{,}4\times10^5\,\mathrm{J/kg}$ và $c=4180\,\mathrm{J/(kg\cdot K)}$.
\choice
{\True $1694400\,\mathrm{J}$}
{$136\times10^4\,\mathrm{J}$}
{$334400\,\mathrm{J}$}
{$1894400\,\mathrm{J}$}
\loigiai{Nhiệt lượng tổng cộng: $Q=m\lambda+mc\Delta T=4\cdot3{,}4\times10^5+4\cdot4180\cdot20=1694400\,\mathrm{J}$.}
\end{ex}

\begin{ex}
Người ta dùng một lò nung điện có công suất $20\,\mathrm{kW}$ để làm nóng chảy hoàn toàn $2\,\mathrm{kg}$ đồng có nhiệt độ ban đầu $30^\circ\mathrm{C}$. Biết chỉ $50\%$ năng lượng tiêu thụ của lò được dùng vào việc làm đồng nóng lên và nóng chảy hoàn toàn ở nhiệt độ $1083^\circ\mathrm{C}$. Cho $c=380\,\mathrm{J/(kg\cdot K)}$, $\lambda=1{,}8\times10^5\,\mathrm{J/kg}$. Thời gian cần thiết khoảng
\choice
{$58\,\mathrm{s}$}
{\True $116\,\mathrm{s}$}
{$90\,\mathrm{s}$}
{$30\,\mathrm{s}$}
\loigiai{
$\Delta T=1083-30=1053\,\mathrm{K}$. Khi đó
\[
Q=mc\Delta T+m\lambda
=2\cdot380\cdot1053+2\cdot1{,}8\times10^5
=1160280\,\mathrm{J}.
\]
Do hiệu suất $\eta=50\%$:
\[
\tau=\frac{Q}{\eta\mathcal{P}}
=\frac{1160280}{0{,}5\cdot20000}
\approx116\,\mathrm{s}.
\]
}
\end{ex}

\begin{center}
\textbf{ĐÁP ÁN}
\end{center}

\[
\boxed{1C\;2D\;3C\;4A\;5A\;6C\;7C\;8B\;9C\;10A\;11B\;12A\;13B\;14B\;15C\;16A\;17A\;18B}
\]
